% Hình 3.3: Quy trình nhận diện và xếp hạng Transition tự động từ đặc tả OpenAPI
% Thiết kế đơn sắc (Black & White / Grayscale) cho in ấn học thuật, mở rộng kích thước ô quyết định chống tràn chữ.
\begin{tikzpicture}[
  >=Stealth,
  scale=0.80, transform shape,
  block/.style={rectangle, draw=black, fill=white, thick, minimum width=3.0cm, minimum height=1.0cm, align=center, rounded corners=3pt, font=\small},
  keyblock/.style={rectangle, draw=black, fill=gray!12, thick, minimum width=3.4cm, minimum height=1.2cm, align=center, rounded corners=3pt, font=\small},
  outnode/.style={rectangle, draw=black, fill=white, thick, minimum width=3.8cm, minimum height=0.85cm, align=center, rounded corners=2pt, font=\footnotesize},
  groupnode/.style={draw=black!60, dashed, fill=gray!5, rounded corners=5pt, inner sep=8pt}
]

  % 1. Đặc tả đầu vào
  \node[block] (spec) at (0, 0) {Đặc tả OpenAPI\\/ Swagger Spec};
  
  % 2. Tách đặc trưng
  \node[block] (tokenizer) at (4.0, 0) {Bộ trích xuất đặc trưng\\(Tokenization \& Cleaning)\\[2pt]
    \scriptsize \texttt{p:token}, \texttt{f:token}, \texttt{v:verb}};

  % 3. Hai nhánh mô hình
  \node[keyblock] (nb)  at (8.6, 1.5)  {\textbf{Mô hình Naive Bayes}\\(huấn luyện trên 10.381\\thao tác ghi của 230 API)};
  \node[block]    (lex) at (8.6, -1.5) {Bộ từ vựng Lexicon\\(53 token đường dẫn,\\45 token tham số)};

  % 4. Khối chấm điểm & nhận diện cấu trúc
  \node[keyblock] (engine) at (13.2, 0) {\textbf{Động cơ chấm điểm}\\và nhận diện cấu trúc};

  % 5. Bốn quyết định đầu ra
  \node[outnode] (out_op)     at (18.2, 1.65)  {\textbf{1. Thao tác chuyển trạng thái} ($\mathrm{score} > 0$)};
  \node[outnode] (out_role)   at (18.2, 0.55)  {\textbf{2. Trường vai trò}};
  \node[outnode] (out_victim) at (18.2, -0.55) {\textbf{3. Trường chỉ tài khoản nạn nhân}};
  \node[outnode] (out_val)    at (18.2, -1.65) {\textbf{4. Giá trị hạ quyền}};

  % Nhóm quyết định
  \begin{scope}[on background layer]
    \node[groupnode, fit=(out_op) (out_role) (out_victim) (out_val), 
          label={[font=\bfseries\footnotesize, text=black]above:Quyết định sinh ca kiểm thử}] (out_group) {};
  \end{scope}

  % Đường nối
  \draw[->, thick] (spec) -- (tokenizer);
  \draw[->, thick] (tokenizer) |- (nb);
  \draw[->, thick, dashed] (tokenizer) |- (lex);
  \draw[->, very thick] (nb) -| (engine);
  \draw[->, thick, dashed] (lex) -| (engine);

  \draw[->, thick] (engine.east) -- (out_op.west);
  \draw[->, thick] (engine.east) -- (out_role.west);
  \draw[->, thick] (engine.east) -- (out_victim.west);
  \draw[->, thick] (engine.east) -- (out_val.west);

\end{tikzpicture}
